%% scholia-manual.tex
%% User Manual for the scholia LaTeX package
%% Copyright 2025-2026 Esther Poniatowski
%% LPPL 1.3c

\documentclass[11pt,a4paper]{article}

% Path setup for local development
% When installed via CTAN/TDS, replace with: \usepackage{scholia}
\def\CommonPath{../..}
\usepackage{\CommonPath/scholia/scholia}

% Additional packages for documentation
\usepackage{hyperref}
\usepackage{listings}
\usepackage{multicol}
\usepackage{booktabs}
\usepackage{fancyvrb}
\usepackage{xspace}

% Hyperref configuration
\hypersetup{
  colorlinks=true,
  linkcolor=callout-method,
  urlcolor=callout-definition,
  pdftitle={scholia - User Manual},
  pdfauthor={Esther Poniatowski},
}

% Custom commands for the manual
\newcommand{\pkg}[1]{\textsf{#1}\xspace}
\newcommand{\cmd}[1]{\texttt{\textbackslash#1}}
\newcommand{\env}[1]{\texttt{#1}}
\newcommand{\opt}[1]{\texttt{#1}}
\newcommand{\scholia}{\textsc{scholia}\xspace}

% Code listing style
\lstset{
  language=[LaTeX]TeX,
  basicstyle=\ttfamily\small,
  keywordstyle=\color{callout-method}\bfseries,
  commentstyle=\color{gray},
  backgroundcolor=\color{gray!10},
  frame=single,
  framerule=0pt,
  xleftmargin=1em,
  xrightmargin=1em,
  breaklines=true,
  showstringspaces=false,
}

% Document metadata
\title{\Huge\bfseries\textsc{scholia}\\[0.5em]
  \Large Structured Visual Components for \LaTeX{} Documents}
\author{Esther Poniatowski\\
  \texttt{esther.poniatowski@ens.psl.eu}}
\date{Version 0.1.0\\January 2026}

\begin{document}

\maketitle

\begin{abstract}
\scholia is a \LaTeX{} package for documents requiring consistent visual structure.
It provides callout boxes and semantic labels for highlighting content,
and numbering systems to organize steps and components. For educational use,
it adds exercise environments, grading tools, and worksheet layouts. The package
supports article and Beamer modes with French/English localization.
\end{abstract}

\tableofcontents
\newpage

%% ===========================================================================
\section{Overview}
%% ===========================================================================

\scholia provides environments and commands for documents requiring structured
visual components. Typical applications include educational materials, technical
tutorials, reports, and academic projects.

\subsection{Key Features}

The package offers six main feature categories, listed from most general
(applicable to any structured document) to most specific (educational content):

\begin{tabular}{@{}lp{9cm}@{}}
\toprule
\textbf{Feature} & \textbf{Description} \\
\midrule
Callouts & 13 styled box types (definitions, theorems, warnings, tips, etc.) \\
Semantic labels & Inline styled boxes for content categorization \\
Numbering & Steps/substeps and questions/subquestions with hanging indentation \\
Typed Pages & Standalone page layouts with banners and relative pagination \\
Exercises & 5 environment types with difficulty indicators and solution control \\
Grading & Rubric tables and point allocation for assessments \\
\bottomrule
\end{tabular}

\subsection{Document Modes}

The package is designed to maintain consistent styling across multiple formats. This allows
presenting the same content in both printed documents and slide presentations.

Two modes are supported:

\begin{tabular}{@{}lp{4cm}p{6cm}@{}}
\toprule
\textbf{Mode} & \textbf{Purpose} & \textbf{Features} \\
\midrule
\textbf{Article} (default) & Printed materials, handouts, worksheets &
  All features enabled. \\
\textbf{Beamer} & Presentations, oral sessions &
  Typed pages disabled (page-based constructs incompatible with slides).
  Semantic labels available. Progression module loaded for slide navigation. \\
\bottomrule
\end{tabular}

\subsection{Pre-Defined Styles and Customization}

\scholia offers two levels of usage:

\begin{itemize}
  \item \textbf{Ready-to-use environments} --- Variants in each feature, with pre-defined colors, icons,
        and styling. No configuration required. (13 callout types, 9 semantic labels, 5 exercise categories,
        10 typed page layouts)
  \item \textbf{Customization framework} --- Generic commands allow creating new variants in each
  feature. All visual elements can be redefined (colors, icons, spacing, borders) using the same underlying infrastructure.
\end{itemize}

%% ===========================================================================
\section{Installation}
%% ===========================================================================

\subsection{CTAN Installation}

Once available on CTAN, install via the standard TeX distribution tools:

\begin{lstlisting}
# TeX Live
tlmgr install scholia

# MiKTeX
miktex-console --install=scholia
\end{lstlisting}

\subsection{Manual Installation}

Place the \texttt{scholia/} directory in a local texmf tree:

\begin{lstlisting}
~/texmf/tex/latex/scholia/
\end{lstlisting}

Then refresh the filename database:

\begin{lstlisting}
texhash ~/texmf
\end{lstlisting}

%% ===========================================================================
\section{Loading the Package}
%% ===========================================================================

\subsection{Basic Usage}

\scholia is loaded in the preamble with \cmd{usepackage}. The package detects
the document class automatically: with \texttt{article}, \texttt{report}, or
\texttt{book}, it operates in article mode; with \texttt{beamer}, it switches
to presentation mode.

\medskip
\textbf{Standard article with English strings (default):}

\begin{lstlisting}
\documentclass{article}
\usepackage{scholia}
\end{lstlisting}

\textbf{Article with French strings:}

\begin{lstlisting}
\documentclass{article}
\usepackage[french]{scholia}
\end{lstlisting}

\textbf{Beamer presentation:}

The \opt{beamer} option can be specified explicitly, though it is auto-detected
when using the \texttt{beamer} document class.

\begin{lstlisting}
\documentclass{beamer}
\usepackage[beamer]{scholia}
\end{lstlisting}

\subsection{Package Options}

\begin{tabular}{@{}lp{10cm}@{}}
\toprule
\textbf{Option} & \textbf{Description} \\
\midrule
\opt{english} & Use English strings (default) \\
\opt{french} & Use French strings \\
\opt{beamer} & Enable Beamer mode (auto-detected for \texttt{beamer} class) \\
\opt{noicons} & Disable icon loading for text-only output \\
\opt{minimal} & Minimal dependencies; implies \opt{noicons} \\
\opt{draft} & Draft mode (reserved for future use) \\
\bottomrule
\end{tabular}

\subsection{Localization}

All user-facing strings (labels, titles, prompts) are defined in localization files.
The \opt{french} and \opt{english} options select the appropriate string set.

\begin{tabular}{@{}lll@{}}
\toprule
\textbf{Element} & \textbf{French} & \textbf{English} \\
\midrule
Exercise label & Exercice & Exercise \\
Definition & D\'efinition & Definition \\
Step & \'Etape & Step \\
Total & Total & Total \\
\bottomrule
\end{tabular}

\medskip
To override specific strings, redefine the internal macros after loading \scholia:

\begin{lstlisting}
\usepackage{scholia}
\def\scholia@str@exo@label{Problem}
\end{lstlisting}

\subsection{Icons}

Callouts and exercises display icons by default, loaded from PDF files in the
\texttt{icons/pdf/} directory. The \opt{noicons} option disables icon loading,
which is useful for:

\begin{itemize}
  \item Faster compilation during drafting
  \item Reduced output file size
  \item Plain-text fallback scenarios
\end{itemize}

%% ===========================================================================
\section{Quick Start}
%% ===========================================================================

The following example demonstrates \scholia in an article document with a definition
callout and a guided exercise. The exercise uses numbered steps and a conditional
answer blank.

\begin{lstlisting}
\documentclass{article}
\usepackage[english]{scholia}

\begin{document}

\begin{definitionbox}[Derivative]
The derivative of $f$ at $x$ is the limit
$f'(x) = \lim_{h \to 0} \frac{f(x+h) - f(x)}{h}$
when this limit exists.
\end{definitionbox}

\begin{exguided}[Computing Derivatives]
\exstep{Find $f'(x)$ for $f(x) = x^2$.}

Apply the definition: $\lim_{h \to 0} \frac{(x+h)^2 - x^2}{h}$.

The answer is \solutionBlank{2cm}{$2x$}.
\end{exguided}

\end{document}
\end{lstlisting}

%% ===========================================================================
\section{Features}
%% ===========================================================================

This section documents each feature category in detail.

%% ---------------------------------------------------------------------------
\subsection{Callouts}
%% ---------------------------------------------------------------------------

Callouts are styled boxes for highlighting specific content types. Each callout
type provides two environment variants:

\begin{itemize}
  \item \env{<type>box} --- Full box with optional title and body content
  \item \env{<type>mark} --- Title-only bar (no body area)
\end{itemize}

\subsubsection{Available Types}

\begin{tabular}{@{}llp{6cm}@{}}
\toprule
\textbf{Environment} & \textbf{Color} & \textbf{Typical Use} \\
\midrule
\env{definitionbox} & Cyan & Mathematical definitions \\
\env{theorembox} & Red & Theorems and propositions \\
\env{proofbox} & Gray & Mathematical proofs \\
\env{warningbox} & Orange & Warnings, common mistakes \\
\env{errorbox} & Dark red & What NOT to do \\
\env{tipbox} & Yellow & Tips and strategies \\
\env{methodbox} & Blue & Step-by-step methods \\
\env{notationbox} & Brown-gray & Notation conventions \\
\env{objectivebox} & Green & Learning objectives \\
\env{summarybox} & Light gray & Summaries and recaps \\
\env{remarkbox} & Purple & Observations \\
\env{questionbox} & Blue & Questions for reflection \\
\env{instructionbox} & Gray & Instructions and directives \\
\bottomrule
\end{tabular}

\subsubsection{Usage Examples}

\paragraph{Box with Title}\leavevmode

\begin{lstlisting}
\begin{definitionbox}[Some title]
Some content here.
\end{definitionbox}
\end{lstlisting}

\begin{definitionbox}[Some title]
Some content here.
\end{definitionbox}

\paragraph{Box without Title}\leavevmode

\begin{lstlisting}
\begin{tipbox}
Some content here.
\end{tipbox}
\end{lstlisting}

\begin{tipbox}
Some content here.
\end{tipbox}

\paragraph{Title-Only Variant}\leavevmode

The \texttt{*mark} variants are commands (not environments) for title-only banners:

\begin{lstlisting}
\theoremmark{Some text}
\end{lstlisting}

\theoremmark{Some text}

The title-only variant accepts an optional first argument for \pkg{tcolorbox} options:

\begin{lstlisting}
\theoremmark[centered]{Centered text}
\end{lstlisting}

\subsubsection{Customization}

Each callout type is defined by a color. To create a new callout type:

\begin{enumerate}
  \item Define a color with \cmd{definecolor}
  \item Create the environment using \texttt{callout layout} (for boxes) or
        \texttt{callout titleonly} (for title-only variants)
\end{enumerate}

\begin{lstlisting}
\definecolor{callout-custom}{RGB}{100,150,200}
\newtcolorbox{custombox}[1][]{%
  callout layout=callout-custom,
  title={#1},
}
\end{lstlisting}

The styles are based on \pkg{tcolorbox}; any tcolorbox option can be added to
further customize spacing, borders, or backgrounds.

%% ---------------------------------------------------------------------------
\subsection{Semantic Labels}
%% ---------------------------------------------------------------------------

Semantic labels are inline styled boxes for categorizing and marking content.
They provide visual consistency when referring to concepts, methods,
or other resources. Available only in article mode.

\subsubsection{Generic Frame}

The \cmd{inlineframebox} command creates a colored inline box:

\begin{lstlisting}
See the \inlineframebox[cyan]{theory section} for details.
\end{lstlisting}

With navigation arrow (starred variant):

\begin{lstlisting}
Continue to \inlineframebox[green]*{exercise}.
\end{lstlisting}

\subsubsection{Semantic Inline Labels}

Pre-defined commands for content categorization:

\begin{tabular}{@{}ll@{}}
\toprule
\textbf{Command} & \textbf{Color/Purpose} \\
\midrule
\cmd{conceptslabel\{text\}} & Cyan (theory, definitions) \\
\cmd{exerciseslabel\{text\}} & Orange (practice activities) \\
\cmd{solutionslabel\{text\}} & Green (answers, corrections) \\
\cmd{methodlabel\{text\}} & Blue (methods, procedures) \\
\cmd{assessmentlabel\{text\}} & Red (formal evaluations) \\
\cmd{researchlabel\{text\}} & Reddish-orange (exploration) \\
\cmd{selfassessmentlabel\{text\}} & Purple (self-evaluation) \\
\cmd{keypointslabel\{text\}} & Gold (key highlights) \\
\cmd{summarylabel\{text\}} & Gray-brown (overviews) \\
\bottomrule
\end{tabular}

All labels accept a starred variant for the navigation arrow:

\begin{lstlisting}
Refer to the \notionlabel*{definition of continuity}.
\end{lstlisting}

%% ---------------------------------------------------------------------------
\subsection{Numbering Systems}
%% ---------------------------------------------------------------------------

\scholia provides two numbering systems for structured content, distinguished by
their visual marker style. Each system offers both environment and inline command
variants.

\subsubsection{Squared Items (Steps)}

Visual style: \texttt{[1]}, \texttt{[2]}, \texttt{[3]}... with lettered sub-items
\texttt{(a)}, \texttt{(b)}, \texttt{(c)}...

The generic environments are \texttt{squareditem} and \texttt{squaredsubitem}.
Semantic aliases \texttt{step} and \texttt{substep} are provided for procedural
content. The subitem counter resets automatically when a new item begins.

\paragraph{Environment Syntax}\leavevmode

\begin{lstlisting}
\begin{step}[Preparation]
Gather all necessary materials.
\begin{substep}
Check equipment.
\end{substep}
\begin{substep}
Prepare workspace.
\end{substep}
\end{step}
\end{lstlisting}

\begin{step}[Preparation]
Gather all necessary materials.
\begin{substep}
Check equipment.
\end{substep}
\begin{substep}
Prepare workspace.
\end{substep}
\end{step}

\paragraph{Inline Syntax}

For compact documents, use \cmd{squaremark} and \cmd{squaresubmark} (or aliases
\cmd{s} and \cmd{subs}):

\begin{lstlisting}
\s Gather materials.
\subs Check equipment.
\subs Prepare workspace.
\end{lstlisting}

\subsubsection{Circled Items (Questions)}

Visual style: \textcircled{1}, \textcircled{2}, \textcircled{3}... with circled
sub-items \textcircled{\scriptsize a}, \textcircled{\scriptsize b}, \textcircled{\scriptsize c}...

The generic environments are \texttt{circleditem} and \texttt{circledsubitem}.
Semantic aliases \texttt{question} and \texttt{subquestion} are provided for
assessment content. The subitem counter resets automatically when a new item begins.

\paragraph{Environment Syntax}\leavevmode

\begin{lstlisting}
\begin{question}
Solve $x^2 - 4 = 0$.
\begin{subquestion}
Factor the left side.
\end{subquestion}
\begin{subquestion}
Find all solutions.
\end{subquestion}
\end{question}
\end{lstlisting}

\begin{question}
Solve $x^2 - 4 = 0$.
\begin{subquestion}
Factor the left side.
\end{subquestion}
\begin{subquestion}
Find all solutions.
\end{subquestion}
\end{question}

\paragraph{Inline Syntax}

For compact lists, use \cmd{circmark} and \cmd{circsubmark} (or aliases \cmd{q} and \cmd{subq}):

\begin{lstlisting}
\q What is the capital of France?
\subq Name two rivers.
\subq Name the largest city.
\end{lstlisting}

Both syntaxes share the same counters and produce identical visual output.

\subsubsection{Counter Reset}

\begin{tabular}{@{}lp{8cm}@{}}
\toprule
\textbf{Counter} & \textbf{Reset Behavior} \\
\midrule
\texttt{scholia@squarednum} & Manual reset only \\
\texttt{scholia@squaredsubnum} & Resets when \texttt{scholia@squarednum} increments \\
\texttt{scholia@circlednum} & Resets when entering an exercise environment \\
\texttt{scholia@circledsubnum} & Resets when \texttt{scholia@circlednum} increments \\
\bottomrule
\end{tabular}

%% ---------------------------------------------------------------------------
\subsection{Typed Pages}
%% ---------------------------------------------------------------------------

Typed pages are standalone page layouts with styled banners and relative page
numbering. They are available only in article mode (disabled in Beamer).

\subsubsection{Page Types}

\begin{tabular}{@{}llp{5cm}@{}}
\toprule
\textbf{Type Key} & \textbf{Color} & \textbf{Purpose} \\
\midrule
\texttt{notion} & Cyan & Concepts and theory \\
\texttt{exercise} & Orange & Exercise sheets \\
\texttt{solution} & Green & Solutions \\
\texttt{autocheck} & Purple & Self-assessment \\
\texttt{method} & Blue & Methods and procedures \\
\texttt{assessment} & Red & Formal assessments \\
\texttt{summary} & Gray-brown & Chapter summaries \\
\texttt{keypoint} & Gold & Key points to memorize \\
\texttt{research} & Reddish-orange & Research activities \\
\bottomrule
\end{tabular}

\medskip
\textbf{Note:} The \texttt{presentation} type exists for Beamer integration but
is handled separately through the progression module in Beamer mode.

\subsubsection{Creating a Typed Page}

The \cmd{newTypedPage} command initializes a new typed page:

\begin{lstlisting}
\newTypedPage{3e-func-N1}{notion}{Functions}{Images and Antecedents}
\end{lstlisting}

Arguments:
\begin{enumerate}
  \item Reference identifier (for cross-referencing)
  \item Type key (from table above)
  \item Chapter title (displayed in gray)
  \item Specific topic (displayed in bold)
\end{enumerate}

This creates a banner with the appropriate icon, color, and titles, and sets up
relative page numbering within the typed page section.

\subsubsection{Full-Page Environments}

For special emphasis, full-page framed environments wrap content in a distinctive
border that spans the entire page. These are breakable across multiple pages.

\paragraph{Generic Full-Page Environment}

The \env{typedpagefull} environment is the generic mechanism for creating
full-page framed content with any page type:

\begin{lstlisting}
\begin{typedpagefull}{ref}{type}{Chapter}{Title}
Content here...
\end{typedpagefull}
\end{lstlisting}

An optional first argument controls styling:

\begin{lstlisting}
\begin{typedpagefull}[norules]{ref}{keypoint}{Chapter}{Title}
Content with header/footer rules suppressed...
\end{typedpagefull}
\end{lstlisting}

\begin{tabular}{@{}lp{8cm}@{}}
\toprule
\textbf{Option} & \textbf{Effect} \\
\midrule
\texttt{norules} & Suppress header/footer separator lines (content preserved) \\
\bottomrule
\end{tabular}

\paragraph{Convenience Wrappers}

Pre-defined environments wrap \env{typedpagefull} for common use cases:

\begin{tabular}{@{}llp{5cm}@{}}
\toprule
\textbf{Environment} & \textbf{Type} & \textbf{Options} \\
\midrule
\env{summarypage} & \texttt{summary} & (default styling) \\
\env{keypointspage} & \texttt{keypoint} & \texttt{norules} \\
\bottomrule
\end{tabular}

\medskip
Example usage:

\begin{lstlisting}
\begin{summarypage}{ref}{Chapter}{Title}
Summary content here...
\end{summarypage}

\begin{keypointpage}{ref}{Chapter}{Title}
Key points here (no header/footer rules)...
\end{keypointpage}
\end{lstlisting}

Arguments: reference identifier, chapter title, specific topic.

\subsubsection{Document Metadata}

Typed pages display author and year information in the footer. Define these in the
preamble:

\begin{lstlisting}
\renewcommand{\Author}{Author Name}
\renewcommand{\Year}{2025--2026}
\end{lstlisting}

\subsubsection{Cross-Referencing}

Reference typed pages using the automatically created labels:

\begin{lstlisting}
See \typedPageRef{3e-func-N1} for the definition.
\end{lstlisting}

%% ---------------------------------------------------------------------------
\subsection{Exercises}
%% ---------------------------------------------------------------------------

\scholia provides five exercise environment types. These types are \emph{visual categories}
that indicate the pedagogical context; the actual content within each environment is
entirely flexible.

\subsubsection{Exercise Types}

\begin{tabular}{@{}lp{8cm}@{}}
\toprule
\textbf{Environment} & \textbf{Visual Indication} \\
\midrule
\env{exguided} & Guided icon; suggests step-by-step assistance \\
\env{exsemiguided} & Semi-guided icon; partial hints expected \\
\env{exauto} & Difficulty stars (1--3); independent work \\
\env{exmodel} & Modeling icon; open-ended problem-solving \\
\env{excorrection} & Correction styling; for displaying solutions \\
\bottomrule
\end{tabular}

\subsubsection{Exercise Environments}

All exercise environments share a common structure: an optional title in brackets,
followed by the exercise content. The \env{exauto} environment additionally
requires a mandatory difficulty level (1--3) displayed as stars.

\paragraph{Standard exercise with title}\leavevmode

\begin{lstlisting}
\begin{exguided}[Solving Quadratic Equations]
\exstep{Identify coefficients $a$, $b$, $c$.}
\exstep{Compute $\Delta = b^2 - 4ac$.}
\exstep{Apply the quadratic formula.}
\end{exguided}
\end{lstlisting}

\begin{exguided}[Solving Quadratic Equations]
\exstep{Identify coefficients $a$, $b$, $c$.}
\exstep{Compute $\Delta = b^2 - 4ac$.}
\exstep{Apply the quadratic formula.}
\end{exguided}

\paragraph{Autonomous exercise with difficulty level}\leavevmode

\begin{lstlisting}
\begin{exauto}[Integration Practice][2]
Compute $\int x^2 e^x \, dx$.
\end{exauto}
\end{lstlisting}

\begin{exauto}[Integration Practice][2]
Compute $\int x^2 e^x \, dx$.
\end{exauto}

\subsubsection{Two-Column Layout}

The \env{extwocol} environment displays content in two columns with a
dotted separator, useful for compact exercise sheets:

\begin{lstlisting}
\begin{extwocol}
\begin{exauto}[][1]
Simplify $\frac{x^2 - 1}{x - 1}$.
\end{exauto}

\begin{exauto}[][1]
Factor $x^2 + 5x + 6$.
\end{exauto}
\end{extwocol}
\end{lstlisting}

\subsubsection{Blanks and Hidden Solutions}

\scholia supports dual compilation: one version with answers hidden (for students)
and one with answers visible (for teachers or answer keys). Blank commands insert
placeholders that reveal their content only when solutions are enabled.

\paragraph{Blank Commands}\leavevmode

\begin{tabular}{@{}lp{7cm}@{}}
\toprule
\textbf{Command} & \textbf{Description} \\
\midrule
\cmd{solutionBlank\{width\}\{answer\}} & Underlined blank or answer \\
\cmd{solutionDotsBlank\{n\}\{answer\}} & Dotted blank (n dots) or answer \\
\cmd{fillBlank\{width\}} & Always show underlined blank \\
\cmd{blankLine\{width\}} & Simple underline \\
\bottomrule
\end{tabular}

\medskip
Example usage:

\begin{lstlisting}
The derivative is \solutionBlank{2cm}{$2x$}.
\end{lstlisting}

\paragraph{Enabling Solutions}

By default, solutions are hidden. Two methods enable solution display:

\textbf{Method 1: In the preamble}

\begin{lstlisting}
\usepackage{scholia}
\solutionstrue
\end{lstlisting}

\textbf{Method 2: Via command line}

\begin{lstlisting}[language=bash]
pdflatex "\def\ForceSolutions{}\input{document}"
\end{lstlisting}

This allows generating both student and teacher versions from the same source file.

\subsubsection{Exercise Commands}

\begin{tabular}{@{}lp{7cm}@{}}
\toprule
\textbf{Command} & \textbf{Description} \\
\midrule
\cmd{exstep[title]} & Numbered step (a, b, c, ...) \\
\cmd{resetExerciseSteps} & Reset step counter to zero \\
\cmd{highlight\{text\}} & Subtle gray background \\
\cmd{questionStatement\{text\}} & Highlighted question in corrections \\
\bottomrule
\end{tabular}

\subsubsection{Counter Management}

The exercise counter (\texttt{exonum}) increments automatically with each exercise
environment. The step counter (\texttt{exoitem}) resets at each \cmd{resetExerciseSteps}
call. The question counter (\texttt{scholia@circlednum}) resets automatically when entering
any exercise environment.

%% ---------------------------------------------------------------------------
\subsection{Grading}
%% ---------------------------------------------------------------------------

The grading module provides tools for creating assessment rubrics and point
allocation displays. These are intended for teacher documents (answer keys,
grading sheets) rather than student-facing materials.

\subsubsection{Grading Box}

A simple point allocation display:

\begin{lstlisting}
\begin{gradebox}[Point Distribution]
\gradeitem{Question 1}{2 pts}
\gradeitem{Question 2}{3 pts}
\gradetotal{5 pts}
\end{gradebox}
\end{lstlisting}

\begin{gradebox}[Point Distribution]
\gradeitem{Question 1}{2 pts}
\gradeitem{Question 2}{3 pts}
\gradetotal{5 pts}
\end{gradebox}

\subsubsection{Grading Table}

For detailed rubrics with criteria descriptions:

\begin{lstlisting}
\begin{gradetable}{Assessment Rubric}
\graderow{Clarity}{Clear presentation}{2 pts}
\graderow{Accuracy}{Correct reasoning}{3 pts}
\graderowtotal{5 pts}
\end{gradetable}
\end{lstlisting}

\subsubsection{Inline Point Indicators}

Point values can be added inline within exercise text:

\begin{lstlisting}
\begin{question}
Calculate $f(3)$. \gradeinline{2 pts}
\end{question}
\end{lstlisting}

\subsubsection{Automatic Point Tracking}

The grading module includes counters for automatic point summation:

\begin{lstlisting}
\resetgradepoints
\gradeitemTracked{Q1}{Description}{2}
\gradeitemTracked{Q2}{Description}{3}
Total: \showgradetotal
\end{lstlisting}

%% ===========================================================================
\section{License}
%% ===========================================================================

\scholia is distributed under the \LaTeX{} Project Public License (LPPL)
version 1.3c or later.

\begin{quote}
This work may be distributed and/or modified under the conditions of
the \LaTeX{} Project Public License, either version 1.3c of this license
or (at your option) any later version.
\end{quote}

The latest version of this license is available at:\\
\url{https://www.latex-project.org/lppl.txt}

%% ===========================================================================
\section{Acknowledgments}
%% ===========================================================================

\scholia uses the following packages:
\pkg{tcolorbox}, \pkg{tikz}, \pkg{xparse}, \pkg{etoolbox},
\pkg{enumitem}, \pkg{xcolor}, \pkg{geometry}, \pkg{fancyhdr},
and \pkg{titlesec}.

Icons are based on the Lucide icon set (\url{https://lucide.dev}).

\end{document}
